% !TeX program = lualatex
% ================================================================
% Polish Parallel Text Version of BearingsChallenge
%
% Generated by the server using the following settings:
%
%   language            Polish (pol)
%   text direction      left to right
%   font                Noto Serif
%   fallback font       Noto Serif
%   built from          latex/worksheets/geometry/bearings/bearingsChallenge.tex
%   vocabulary key      latex/worksheets/geometry/bearings/bearingsChallenge/L2/pol/bearingsChallenge_polishKey.tex
%   tier 3 vocabulary   green   RGB 0 128 0
%   English text        black   RGB 0 0 0
%   Polish text         purple  RGB 112 48 160
%   layout macros       version 3
%
% Please don't edit any information in this comment block - the server
% compares it against the current settings and rebuilds this file when
% they differ, so changes here would be overwritten. However, feel free
% to make updates to the LaTeX below if you want to edit it.
% ================================================================

\documentclass{article}
% ================================================================
% Copied from the English sheet this was translated from, so that
% anything it sets up for itself still works here
% ================================================================
\usepackage{gensymb}
\usepackage{amsfonts}
\usepackage{graphicx} % Required for inserting images
\usepackage{tabularx}

\title{Bearings Challenge}

\date{}

% Characters this language's font has not got are borrowed from another,
% rather than being silently dropped - see docs/translated-sheet-latex.md
\directlua{%
  if luaotfload and luaotfload.add_fallback then
    luaotfload.add_fallback("eallatinfallback",
      { "Noto Serif:mode=harf;" })
  else
    tex.error("this luaotfload is too old to register a fallback font")
  end
}
\usepackage[english]{babel}
\babelprovide[import]{polish}
\babelfont[polish]{rm}[Renderer=Harfbuzz, BoldFont={Noto Serif}, ItalicFont={Noto Serif}, BoldItalicFont={Noto Serif}, RawFeature={fallback=eallatinfallback}]{Noto Serif}
\babelfont[polish]{sf}[Renderer=Harfbuzz, BoldFont={Noto Serif}, ItalicFont={Noto Serif}, BoldItalicFont={Noto Serif}, RawFeature={fallback=eallatinfallback}]{Noto Serif}
\newcommand{\ealtext}[1]{\foreignlanguage{polish}{#1}}
\newcommand{\ealtextblock}[1]{{\raggedright\foreignlanguage{polish}{#1}\par}}
% ================================================================
% EAL layout helpers
% ================================================================
\usepackage{xcolor}

\definecolor{ealkeycolour}{RGB}{0,128,0}
\definecolor{ealenglish}{RGB}{0,0,0}
\definecolor{ealtwocolour}{RGB}{112,48,160}

\newsavebox{\ealboxen}
\newsavebox{\ealboxtr}
\newlength{\ealglosswd}
\newlength{\ealglossraise}
\setlength{\ealglossraise}{1.15em}

% a tier 3 word, in English and in translation alike
\newcommand{\ealkey}[1]{\textcolor{ealkeycolour}{#1}}
\newcommand{\ealkeytr}[1]{\textcolor{ealkeycolour}{\ealtext{#1}}}

% One translated word above one English word. Built from kernel
% primitives, never a tabular, which would break wherever the sheet
% loads array, booktabs or tabularx. The box is as wide as the wider of
% the two, so a long translation pushes its neighbours aside instead of
% overlapping them, and it claims its full height so a table row or a
% TikZ node makes room for it.
\newcommand{\ealgloss}[2]{%
  \sbox{\ealboxen}{\ealkey{#1}}%
  \sbox{\ealboxtr}{\small\textcolor{ealtwocolour}{\ealtext{#2}}}%
  \setlength{\ealglosswd}{\wd\ealboxen}%
  \ifdim\wd\ealboxtr>\ealglosswd\setlength{\ealglosswd}{\wd\ealboxtr}\fi
  \leavevmode
  \makebox[\ealglosswd][c]{%
    \makebox[0pt][c]{\usebox{\ealboxen}}%
    \makebox[0pt][c]{\raisebox{\ealglossraise}{\usebox{\ealboxtr}}}%
  }%
}

% opens up line spacing around a block of glosses, so the raised
% translations sit clear of the line above
\newenvironment{ealglossed}{\par\linespread{2.1}\selectfont\ignorespaces}{\par}

% a whole translated sentence above its English counterpart. block level,
% so it does not belong inside a tabular cell
\newcommand{\ealpara}[2]{%
  \par\addvspace{0.5\baselineskip}%
  {\color{ealtwocolour}\ealtextblock{#1}}%
  \penalty200
  {\color{ealenglish}#2\par}%
  \addvspace{0.5\baselineskip}%
}

\begin{document}

\maketitle

\ealpara{Bob zna położenie pewnego zakopanego złota; leży ono dokładnie $100$km stąd, na linii dokładnie $123.5 \degree$ zgodnie z ruchem wskazówek zegara od północy.}{Bob knows the location of some buried gold, it lies exactly $100$km away, on a line exactly $123.5 \degree$ clockwise from North.}

\ealpara{Bob ma elektroniczny kompas, który może kierować nim według \ealkeytr{azymutów}, bez błędu; niestety, zaakceptuje tylko 3-cyfrowe \ealkeytr{azymuty}, bez \ealkeytr{ułamków} \ealkeytr{stopni}.}{Bob has an electronic compass that can direct him on \ealkey{bearings}, with no error, unfortunately, it will only accept 3 digit \ealkey{bearings}, with no \ealkey{fractions} of \ealkey{degrees}.}

\ealpara{Aby znaleźć złoto, Bob musi dotrzeć na odległość do $1$m od jego położenia.}{To find the gold, Bob must arrive within $1$m of its location.}

\ealpara{Wprowadzenie}{Introduction}
\section{Introduction}

\ealpara{Załóż, że wszyscy podróżni mogą doskonale mierzyć przebytą odległość.}{Assume all travellers can perfectly measure their distance travelled.}

\begin{enumerate}
    \item \ealpara{\ealkeytr{Skalując} w stosunku $1$cm $:$ $20$km, naszkicuj mapę dla Boba, oznacz \ealkeytr{kąty} i odległości}{Using a \ealkey{scale} of $1$cm $:$ $20$km, sketch a map for Bob, label the \ealkey{angles} and distances}
    \item \ealpara{Ustal, jak Bob może zlokalizować złoto, podróżując po \ealkeytr{azymucie} $180\degree$, a następnie $090\degree$}{Determine how Bob can locate the gold by travelling on a \ealkey{bearing} of $180\degree$ followed by $090\degree$}
    \item \ealpara{Bob może przebyć $20$km dziennie; ile czasu zajmie mu podróż do złota?}{Bob can travel $20$km a day, how long will it take him to travel to the gold?}
    \item \ealpara{Rob natknął się na mapę do złota i zamierza wyznaczyć kurs, aby dotrzeć przed Bobem, podróżując na południowy wschód, a następnie na wschód. Uzbrojony w ten sam model kompasu elektronicznego i podróżując z tą samą prędkością, o ile wcześniej dotrze?}{Rob has stumbled upon the map to the gold, and intends to plot a course to arrive before Bob, travelling Southeast, then East. Armed with the same model of electronic compass, and travelling at the same speed, how much earlier will he arrive?}
    \item \ealpara{Chcesz dotrzeć do skarbu przed Bobem i Robem, ale także możesz podążać tylko za 3-cyfrowymi \ealkeytr{azymutami}. A teraz obaj mają $7$ godzin przewagi! Pokonując $20$km dziennie, czy możesz znaleźć sposób, aby dotrzeć do złota pierwszy? Czy twoja trasa jest najszybsza możliwa?}{You want to beat both Bob and Rob to the treasure, but can also only follow $3$ digit \ealkey{bearings}. And now they both have a $7$ hour head start! Covering $20$km per day can you find a way to arrive at the gold first? is your route the quickest possible?}
\end{enumerate}

\newpage

\ealpara{Wyzwanie I}{Challenge I}
\section{Challenge I}

\ealpara{Wcześniej zakładaliśmy, że podróżni mogą mierzyć swoje odległości z nieograniczoną dokładnością – w rzeczywistości tak nie jest. Na szczęście ty, Bob i Rob pobraliście aplikację, która da sygnał za każdym razem, gdy przebędziesz dokładnie $1$km po linii prostej.}{Previously we assumed that the travellers could measure their distances with unlimited precision - in reality that is not the case. Luckily, you, Bob and Rob have downloaded an app that will ping every time you travel exactly $1$km in a straight line.}

\ealpara{Dlatego wszystkie odcinki podróży muszą być dokładnymi wielokrotnościami $1$km!}{Therefore - all sections of journeys must be exact multiples of $1$km!}

\begin{enumerate}
    \item \ealpara{Ustal, jak daleko od położenia złota znajdziesz się ty, Bob i Rob, jeśli \ealkeytr{zaokrąglisz} każdy odcinek swoich zaplanowanych podróży do najbliższego kilometra. W ten sposób pokaż, że złoto pozostanie zakopane.}{Determine how far off the gold's location you, Bob and Rob will be if you \ealkey{round} each leg of your planned journeys to the nearest kilometre. Thus show that the gold will remain buried.}
    \item \ealpara{Jaki jest wynik podróżowania następującą trasą:}{What is the outcome of travelling the following route:}
    \begin{itemize}
    \item \ealpara{$2$km po \ealkeytr{azymucie} $179\degree$}{$2$km on a \ealkey{bearing} $179\degree$}
    \item \ealpara{$2$km po \ealkeytr{azymucie} $000\degree$}{$2$km on a \ealkey{bearing} of $000\degree$}
    \item \ealpara{$1$km po \ealkeytr{azymucie} $182\degree$}{$1$km on a \ealkey{bearing} of $182\degree$}
    \item \ealpara{$1$km po \ealkeytr{azymucie} $000\degree$}{$1$km on a \ealkey{bearing} of $000\degree$}
    \end{itemize}
    \item \ealpara{Jaki jest wynik podróżowania następującą trasą:}{What is the outcome of travelling the following route:}
    \begin{itemize}
        \item \ealpara{$2$km po \ealkeytr{azymucie} $091\degree$}{$2$km on a \ealkey{bearing} $091\degree$}
        \item \ealpara{$2$km po \ealkeytr{azymucie} $270\degree$}{$2$km on a \ealkey{bearing} of $270\degree$}
        \item \ealpara{$1$km po \ealkeytr{azymucie} $088\degree$}{$1$km on a \ealkey{bearing} of $088\degree$}
        \item \ealpara{$1$km po \ealkeytr{azymucie} $270\degree$}{$1$km on a \ealkey{bearing} of $270\degree$}
    \end{itemize}
    \item \ealpara{Użyj poprzednich dwóch pytań, aby pokazać, że trasa prowadząca do złota jest możliwa – ile czasu zajmie ta trasa?}{Use the previous two questions to show that a route leading to the gold is possible - how long will this route take?}
    \item \ealpara{Pokaż, że można wyznaczyć trasę na odległość do $1$m od dowolnego położenia $(x,y)$}{Show that it is possible to plot a route to within $1$m of any location $(x,y)$}
    \item \ealpara{Pokaż, że można wyznaczyć trasę o co najwyżej $6$ odcinkach na odległość do $1$m od dowolnego położenia $(x,y)$}{Show that it is possible to plot a route with at most $6$ legs to within $1$m of any location $(x,y)$}
\end{enumerate}

\vspace{1em}
\ealpara{Wyzwanie II}{Challenge II}
\section{Challenge II}

\ealpara{Teraz ograniczamy naszą uwagę do tras, które mają tylko dwa odcinki, a więc przyjmują następującą postać:}{Now we restrict our attention to routes that have only two legs, thus take the following form:}

\begin{itemize}
    \item \ealpara{Przebądź $M$km po \ealkeytr{azymucie} $ABC\degree$}{Travel $M$km on a \ealkey{bearing} of $ABC\degree$}
    \item \ealpara{Przebądź $N$km po \ealkeytr{azymucie} $XYZ\degree$}{Travel $N$km on a \ealkey{bearing} of $XYZ\degree$}
\end{itemize}

\ealpara{Gdzie $000\degree \leq ABC, XYZ < 360\degree$ oraz $M,N \in \mathbb{N}$}{Where $000\degree \leq ABC, XYZ < 360\degree$, and $M,N \in \mathbb{N}$}

\begin{enumerate}
    \item \ealpara{Rozważ przypadek, w którym $M>6000$; jakie ograniczenie nakłada to na $XYZ\degree$, aby powrócić na odległość $100$km od miejsca startu?}{Consider the case that $M>6000$, what restriction does this place on $XYZ\degree$ to return within $100$km of the start location?}
    \item \ealpara{Użyj \ealkeytr{twierdzenia cosinusów}, aby wyznaczyć \ealkeytr{rozstęp} możliwych wartości $N$ dla danego $M$, aby koniec podróży znajdował się między $0$ a $100$km od startu}{Use the \ealkey{cosine rule} to determine the \ealkey{range} of $N$ values possible given $M$, to ensure that the end of the journey lies between $0$ and $100$km from the start}
    \item \ealpara{W ten sposób ustal \ealkeytr{ograniczenie górne} liczby możliwych dwuczłonowych podróży, które musimy rozważyć}{Thus establish an \ealkey{upper bound} on the number of possible two leg journeys that we need to consider}
    \item \ealpara{Czy wniosek z poprzedniej sekcji nadal musi być prawdziwy?}{Does the conclusion from the previous section still necessarily hold?}
    \item \ealpara{Użyj poniższej tabeli, aby udoskonalić \ealkeytr{ograniczenie górne} liczby możliwych dwuczłonowych podróży do rozważenia:}{Use the following table to refine the \ealkey{upper bound} of the number of possible two leg journeys to consider:}

\begin{table}[h]
\centering
\begin{tabularx}{\textwidth}{|l|X|X|X|X|X|}
\hline
 & $ABC^{\circ}$ Option count &  $M$ in \ealkey{range} &  $XYZ^{\circ}$ Option count per $M$ value (\ealkey{upper bound}) & $N$ Option count per M value (\ealkey{upper bound}) & Total journey option count for this $M$ \ealkey{range} \\
\hline
$M \le 100$            & & & & & \\
\hline
$100 < M \le 500$      & & & & & \\
\hline
$500 < M \le 2000$     & & & & & \\
\hline
$2000 < M \le 6000$    & & & & & \\
\hline
\end{tabularx}
\end{table}

\item \ealpara{Rozważ podróż:}{Consider the journey:}
\begin{itemize}
    \item \ealpara{Przebądź $N$km po \ealkeytr{azymucie} $XYZ\degree$}{Travel $N$km on a \ealkey{bearing} of $XYZ\degree$}
    \item \ealpara{Przebądź $M$km po \ealkeytr{azymucie} $ABC\degree$}{Travel $M$km on a \ealkey{bearing} of $ABC\degree$}
\end{itemize}
\ealpara{a zatem ustal \ealkeytr{ograniczenie górne} $U$ liczby dwuczłonowych podróży, które kończą się w odległości $100$km od punktu startu, przy $U < 10^{10}$}{and thus establish and \ealkey{upper bound} $U$, on the number of two leg journeys that finish within $100$km of the start point, with $U < 10^{10}$}

\item \ealpara{Rozważ \ealkeytr{pole} w metrach kwadratowych koła o \ealkeytr{promieniu} $100$km, a zatem pokaż, że istnieją pewne położenia, których nie można osiągnąć na odległość $1$m, podróżując dwuczłonowo}{Consider the \ealkey{area} in square meters of a circle \ealkey{radius} $100$km, thus show that there are some locations that cannot be reached within $1$m, following a two leg journey}
\item \ealpara{Użyj komputera, aby ustalić, czy złoto można osiągnąć w dwuczłonowej podróży}{Use a computer to determine if the gold can be reached via a two leg journey}
\item \ealpara{Jeśli złoto można osiągnąć, użyj komputera, aby znaleźć najbliższy punkt do złota, którego nie można osiągnąć; jeśli złota nie można osiągnąć, znajdź punkt najbliższy złota, który można osiągnąć}{If the gold can be reached, use a computer to find the closest point to the gold that cannot be reached; if the gold cannot be reached, find the point closest to the gold that can}

\end{enumerate}

\end{document}